% ---------------------------------------------------------------------------
% Second application: switched power converters.  Every number below is backed
% by a file on aeonlab2024 (see comments).  Figures referenced:
%   fig_conv_novel_g20_s0.pdf   (results/figs_conv)  the four discovered converters, added device in red
%   fig_conv_parents_g20_s0.pdf (results/figs_conv)  their nearest training samples
%   table_merged.tex            (results/conv_mech_verify) -> \label{tab:conv_parent_child}
% ---------------------------------------------------------------------------
\section{Second application: switched power converters}
\label{sec:converters}

\paragraph{Dataset, representation and generator.}
To test whether the loop transfers beyond op-amps we re-targeted it at the
LaMAGIC converter corpus~\cite{lamagic}, a set of small hard-switched
DC--DC topologies built from four two-terminal device types: a phase-$a$
switch $S_a$ (closed for the duty fraction $D$ of the period), a phase-$b$
switch $S_b$ (closed for $1-D$), an inductor and a capacitor, wired between
the rails $V_\mathrm{IN}$, $V_\mathrm{OUT}$, ground and internal nets.
Each record is a device--net hypergraph with three to six devices, a duty
cycle from $\{0.1,0.3,0.5,0.7,0.9\}$, and ngspice labels for $V_\mathrm{out}$
and efficiency. We read the two published shards (3--5 and 6 component),
dropped 1{,}959 records with a missing rail, and split the remaining 196{,}376
graphs 90/5/5 into 176{,}738 / 9{,}818 / 9{,}820 for training, validation and
test. Because LaMAGIC enumerates every duty for every topology, the training
split contains only 37{,}514 structurally distinct circuits (Weisfeiler--Lehman
hash); this distinct set is also the reference against which we score
novelty. The representation is the corpus's own bipartite device--net graph,
with no motif layer: at three to six devices there is nothing to abstract, and
keeping LaMAGIC's formulation keeps the comparison to their generator clean.
Devices and nets occupy eight of the sixteen node classes of the op-amp
encoder (rails are folded into the net they touch), every device--net
incidence is a single undirected edge class since the devices are ideal and
have no polarity, and the only conditioning signal is a one-hot duty label.
We deliberately do \emph{not} condition on $V_\mathrm{out}$ or efficiency:
performance-conditioned generation is the LaMAGIC baseline, whereas our claim
is that the prior supplies the feasible region and the agents plus external
simulation do the discovering. The same discrete diffusion architecture as
the op-amp generator was trained on this graph set for 60 epochs (validation
NLL 65.2, test NLL 64.8). Sampled at $D=0.5$, 71--83\,\% of its outputs pass
the structural validity gate, but none of 100 samples met a 2:1 step-down
specification at a grid duty, and 20--30\,\% are pass-throughs whose output
ignores the switches. The generator therefore knows what a converter looks
like, not what it does, which is exactly the regime the discovery loop is
meant for.

\paragraph{Discovery.}
The loop of Section~\ref{sec:method} runs unchanged; only the three domain
plugs are swapped. The \emph{evaluator} builds an ngspice transient deck at a
stated, conventional operating point (100\,V input, 100\,kHz, 500\,$\mu$H with
0.25\,$\Omega$ winding resistance, 10\,$\mu$F with 30\,m$\Omega$ ESR, 50\,m$\Omega$
switches, 50\,$\Omega$ load) and asks the topology to deliver a 2:1
step-down: $|V_\mathrm{out}-50\,\mathrm{V}|\le 5\,\%$, ripple $\le 5\,\%$, an output
that responds to duty, and efficiency $\ge 90\,\%$, at \emph{some} duty found
by a five-point sweep plus two secant refinements. Efficiency at that point
takes the role that gain played for op-amps, and the four robustness
``seeds'' become four operating corners: nominal, heavy load ($R/2$), light
load ($2R$), and low line ($0.9\,V_\mathrm{IN}$ with L and C 20\,\% low).
The \emph{checker} names faulty terminals (unbound pins, floating nets,
shoot-through switch pairs, a shorted output, no input-to-output path), and
because the genome \emph{is} the device--net graph, every LLM repair or
rescue edit (connect a pin to a net, add an L or C, swap a switch phase,
rewire) is applied in genome space, validity-gated as a set, and re-checked
before any simulation. \emph{Novelty} is one minus the cosine similarity of
WL features to the nearest of the 37{,}514 distinct training circuits. With
these plugs, generation~0 is drawn by sequential Monte Carlo with the
validity gate as its health signal, NSGA-II (population 48, 20 generations)
optimises efficiency and novelty under feasibility domination, variation is
SDEdit re-noising of parents, and the same single LLM controller chooses one
action per generation among VERIFY, RACING, RESCUE, REPAIR and STEER from the
observed population state. Every circuit that ever reaches a Pareto front is
archived and re-scored on all four corners at the end, so the deliverable is
the discovered set, not the final front. Two seeds were run on 24 ngspice
workers; seed 0 spent 1{,}193 corner evaluations and 21 controller calls and
archived 10 corner-robust converters, seed 1 spent 1{,}793 and archived 5.

\paragraph{Evaluation.}
Corner robustness is necessary but not sufficient for a discovery claim, so
we evaluate the archived set in three further steps. First, \emph{exact
overlap with the corpus}: each archived netlist is canonicalised (WL
colouring of the device--net graph with rails fixed) and compared with all
37{,}496 distinct topologies of train, validation and test; for every
non-match we also test whether deleting one or two devices yields a corpus
circuit, which identifies the nearest training sample and the device the
search added. This step removed six archived circuits whose only departure
from a training buck was a floating loop that carries no current, a case the
structural gate admits and the novelty score rewards; we report it as a
limitation and prune such subgraphs before scoring in future runs. Second,
\emph{child-versus-parent re-simulation}: the surviving circuit and its
nearest training sample are re-run under the identical specification at all
four corners, so the effect of the single added device is read directly as a
difference in feasible corners, median and worst-corner efficiency, and
ripple. Third, \emph{mechanism verification} from the waveforms of the same
deck, with a zero-volt probe source in series with every switch and
capacitor: the RMS, peak and mean current of the added device, its
utilisation $U$ (fraction of the period with $|i|>1\,\%$ of the load
current), its peak voltage stress, the per-element dissipated power
$\langle v\,i\rangle$ (conduction and ESR only, since the switches are ideal),
capacitor charge balance, an open-loop load step ($50\to25\to50\,\Omega$),
load and line sweeps, an eight-point component-tolerance grid, and an
equivalent-circuit reduction in which the added device is merged, shorted
or removed and the circuit re-simulated. Table~\ref{tab:conv_parent_child}
collects these measurements for the four converters of seed~0 that are
absent from the corpus and electrically distinct from the textbook buck
(Fig.~\ref{fig:conv_novel}); their nearest training samples are shown in
Fig.~\ref{fig:conv_parents}.

\input{table_merged}   % -> \label{tab:conv_parent_child}

% Optional results paragraph (keep or fold into the discussion):
\paragraph{What was found.}
All fifteen archived converters are buck derivatives; no boost or buck-boost
topology became corner-robust in either seed, and two archived circuits are
the textbook buck itself (train samples \#6935 and \#40128). Of the four
electrically distinct additions, one (a) raises efficiency at every load and
corner and is the best circuit of the run at 99.60\,\%; the reduction test
shows it is a buck with half the inductance and half the winding resistance,
and the gain is entirely the resistance halving. One (c) halves ripple and is
equivalent to a larger output capacitor. One (b) carries no current and is
redundant, a novelty-by-isomorphism that the performance tables expose. One
(d) is a genuine new energy-transfer path: its parent is not a converter at
all, and the added capacitor supplies the freewheel path that makes it meet
the specification at all four corners, at the cost of 6.5\,A and 7.2\,A current
peaks through the capacitor and its switch at a 1\,A load, and 0.4 percentage
points of efficiency relative to the textbook buck. The search thus
recovered one useful trade-off and one new mechanism from a generator that
never produced a feasible converter on its own, and the evaluation chain
correctly separated these from the two additions that only look novel.
